% =====================================================================
% Explainable AI-Based Personalized Learning Path Recommendation System
% for University Students
%
% REVISED VERSION — every reported number is measured and reproducible.
% Source: github.com/abhay27singh/explainable-learning-path-recommender
%
% Title unchanged (assigned by the department).
%
% PAGE BUDGET: this compiles longer than six pages with all seven tables.
% If trimming is needed, remove in this order:
%   1. Table VI (cold-start)      2. Table V (calibration)
%   3. Table VIII (dataset)       4. Table IX (efficiency)
% Never remove Tables I, III or IV — I is the headline, III is the
% contribution, IV is the recommendation evidence.
% =====================================================================

\documentclass[conference]{IEEEtran}
\IEEEoverridecommandlockouts

\usepackage{cite}
\usepackage{amsmath,amssymb,amsfonts}
\usepackage{algorithm}
\usepackage{algpseudocode}
\usepackage{graphicx}
\usepackage{booktabs}
\usepackage{textcomp}
\usepackage{xcolor}
\def\BibTeX{{\rm B\kern-.05em{\sc i\kern-.025em b}\kern-.08em T\kern-.1667em\lower.7ex\hbox{E}\kern-.125emX}}

\begin{document}

\title{Explainable AI-Based Personalized Learning Path\\
Recommendation System for University Students}

\author{
\IEEEauthorblockN{Abhay Pratap Singh}
\IEEEauthorblockA{\textit{Department of Computer Science}\\
\textit{Chandigarh University}\\
Mohali, India\\
23BCS11784@cuchd.in}
\and
\IEEEauthorblockN{Akshay Guleria}
\IEEEauthorblockA{\textit{Department of Computer Science}\\
\textit{Chandigarh University}\\
Mohali, India\\
23BCS10517@cuchd.in}
\and
\IEEEauthorblockN{Uday Chaturvedi}
\IEEEauthorblockA{\textit{Department of Computer Science}\\
\textit{Chandigarh University}\\
Mohali, India\\
23BCS10807@cuchd.in}
\linebreakand
\IEEEauthorblockN{Gursimar Makkar}
\IEEEauthorblockA{\textit{Department of Computer Science}\\
\textit{Chandigarh University}\\
Mohali, India\\
23BCS11455@cuchd.in}
\and
\IEEEauthorblockN{Shikha Kamal}
\IEEEauthorblockA{\textit{Department of Computer Science}\\
\textit{Chandigarh University}\\
Mohali, India\\
shikha.e12552@cumail.in}
}

\maketitle

% =====================================================================
\begin{abstract}
Digital learning platforms in higher education generate large volumes of interaction
data that could support personalised learning-path recommendation, but existing
educational recommenders offer little transparency and are rarely evaluated in a way
that establishes which of their components actually contribute. This paper presents an
explainable learning-path recommendation framework combining knowledge-graph-conditioned
knowledge tracing, prerequisite-constrained planning, and priority-backtracking
explanation, evaluated on the Open University Learning Analytics Dataset (OULAD). The
knowledge tracing model attains AUC-ROC $0.9538 \pm 0.0019$ and expected calibration
error $0.0054$ under five-fold cross-validation with folds grouped by student,
significantly exceeding Deep Knowledge Tracing ($0.9435 \pm 0.0017$, $p = 0.0009$) and a
graph-based baseline ($0.9502 \pm 0.0014$, $p = 0.022$) after Holm correction. A central
finding concerns evaluation rather than architecture. Removing knowledge-graph
propagation changes predictive accuracy by $-0.0001$ ($p = 0.77$), which would ordinarily
be read as the graph contributing nothing. That reading is an artefact of the
measurement: accuracy is computed only at assessment positions, and only 89 of 237
concepts carry assessments. Under a cold-concept protocol that withholds all assessments
for a subset of concepts and scores only those, the same graph improves accuracy by
$+0.0325$ ($p = 0.0003$, Cohen's $d = 5.06$), winning on every fold. Knowledge-graph
propagation therefore contributes nothing where direct supervision is plentiful and
substantially where it is absent, and conventional ablation cannot distinguish the two
cases. Planning is evaluated against the concepts learners actually studied next, with
zero prerequisite violations across 906 recommendation decisions, and explanation quality
is assessed by objective fidelity, sufficiency and stability rather than by subjective
rating. All code, data transformations and results are released publicly.
\end{abstract}

\begin{IEEEkeywords}
Explainable AI, knowledge tracing, knowledge graphs, learning path recommendation,
educational data mining, ablation methodology.
\end{IEEEkeywords}

% =====================================================================
\section{Introduction}

The rapid digitisation of higher education has created unprecedented opportunities for
data-driven personalisation, yet the one-size-fits-all curriculum remains common in most
universities \cite{li2026survey}. Students choosing between electives, or working through
a compressed degree programme, rarely receive guidance matched to their prior
experience, pace of study and intended direction. In cumulative subjects a gap in a
foundational concept can produce persistent difficulty for the remainder of a programme.

Adaptive recommender systems are intended to address this, but two limitations recur.
First, many systems operate as black boxes, offering no account of why a particular
activity was suggested \cite{ma2024survey}. Students and academic advisers need that
account in order to judge whether a suggestion fits their situation; without it, even
well-founded recommendations are disregarded. Second, personalisation work has
concentrated on adaptive feedback delivery rather than on modelling how a learner's
knowledge evolves and constructing coherent sequences over it.

A third limitation is less often discussed and motivates the present work. Systems of
this kind are routinely assembled from several components --- a sequence model, a
structural prior such as a knowledge graph, a planner --- and each component is justified
by an ablation showing that removing it degrades accuracy. We show that this evaluation
can be systematically uninformative. When a structural prior exists precisely to support
items for which direct supervision is scarce, and accuracy is measured only where
supervision is abundant, the ablation measures the component in the one regime where it
cannot matter.

This paper makes the following contributions.

\begin{enumerate}
\item A knowledge-graph-conditioned knowledge tracing model with a learned per-head
temporal decay, attaining AUC-ROC $0.9538 \pm 0.0019$ and expected calibration error
$0.0054$ on OULAD under student-grouped five-fold cross-validation.

\item Evidence that conventional ablation cannot detect the contribution of a knowledge
graph, together with a cold-concept protocol that can. Prerequisite propagation changes
standard accuracy by $-0.0001$ ($p = 0.77$) yet improves accuracy on concepts lacking
direct supervision by $+0.0325$ ($p = 0.0003$, $d = 5.06$).

\item A prerequisite-constrained planner whose pedagogical guarantee is enforced
structurally rather than learned, verified at zero violations across 906 recommendation
decisions evaluated against observed learner behaviour.

\item Three objective measures of explanation quality --- fidelity, sufficiency and
stability --- in place of subjective rating. One of these revealed that generated
explanations accounted for only $0.2\%$ of the predicted benefit they described, a defect
no rating scale would surface.

\item A reproducible pipeline from raw data to every reported figure, together with a
deployed service, released publicly.
\end{enumerate}

The remainder of the paper is organised as follows. Section~\ref{sec:related} reviews
related work; Section~\ref{sec:method} specifies the framework;
Section~\ref{sec:implementation} details the implementation;
Section~\ref{sec:results} reports results; Section~\ref{sec:conclusion} concludes.

% =====================================================================
\section{Related Work}
\label{sec:related}

Educational recommender systems have progressed from undifferentiated delivery towards
adaptive personalisation. Early systems relied on content-based and collaborative
filtering, which suffer from cold-start difficulties, data sparsity and limited
scalability \cite{li2026survey, aher2013combination}. Hybrid models combining several
recommendation strategies improved accuracy and robustness, but concentrated on
predicting learner preference rather than modelling knowledge acquisition or optimising
sequential pathways.

Deep Knowledge Tracing (DKT) \cite{piech2015dkt} models a learner's latent knowledge
state over time using recurrent networks. Unlike Bayesian Knowledge Tracing
\cite{corbett1994knowledge}, which relies on static per-skill parameters, DKT learns
temporal dependencies directly from interaction data. Subsequent work introduced memory
structures \cite{zhang2017dkvmn} and attention \cite{pandey2019sakt, ghosh2020akt},
the latter allowing a model to attend directly to any earlier interaction rather than
propagating information through a recurrent chain \cite{vaswani2017attention}.

Reinforcement learning has been applied to learning-path recommendation by formalising
the problem as a Markov decision process and optimising for long-term learning gain
rather than immediate engagement. Knowledge graphs offer a complementary route to
structured and explainable recommendation by modelling concepts, course units and their
dependencies explicitly \cite{li2026survey}. Graph neural networks have proven effective
for the relational structure such graphs induce.

Several gaps remain. Explainability is frequently added after the fact rather than
treated as a design requirement, and where it is evaluated at all, evaluation is usually
by participant rating --- which measures whether readers found an explanation agreeable,
not whether it is faithful to the model producing it. Temporal dynamics, particularly
forgetting, are often omitted, which compromises the pedagogical validity of suggested
trajectories. Most relevant to this paper, the contribution of structural priors is
customarily established by an ablation on aggregate predictive accuracy, without
examining whether that metric is capable of observing the component in question. We
return to this point in Section~\ref{sec:results}.

% =====================================================================
\section{Methodology}
\label{sec:method}

\subsection{Overview}

The framework has three components: a knowledge tracing model that estimates a learner's
mastery over every concept, a planner that selects the next concept subject to
prerequisite constraints, and an explanation module that justifies the selection by
reference to the prerequisite graph. Explanation is post-hoc: it operates on a decision
already taken. What distinguishes it from a generic attribution method is that it is
grounded in graph structure and in counterfactuals the model can be queried for
directly, and that its quality is measured objectively rather than surveyed.

\subsection{Concept Graph Construction}

OULAD contains no concept annotations, prerequisite relations, or titles for learning
materials. A concept graph is therefore derived rather than read off. A concept is
defined as a (module, week-of-study) pair: the set of material a module presents in a
given week. Week placement uses the \texttt{week\_from} field where present (17.6\% of
materials) and otherwise the median day on which learners accessed the material. On the
1{,}121 materials carrying both signals the two agree at Spearman $\rho = 0.975$, which
we report as validation of the substitute.

Edges arise from three rules. Sequence edges connect consecutive weeks within a module.
Assessment-coverage edges connect the four preceding weeks to the week an assessment
falls in. Mined edges are admitted where learner behaviour shows both consistent ordering
and a measurable effect: for concepts $a, b$ within a module,
%
\begin{equation}
\mathrm{prec}(a,b) = \frac{\left|\{u : \tau_u(a) < \tau_u(b)\}\right|}{\left|\{u : a, b \in h_u\}\right|},
\quad
\mathrm{lift}(a,b) = \frac{P(y_b \mid a \prec b)}{P(y_b)},
\label{eq:mining}
\end{equation}
%
where $\tau_u(\cdot)$ is learner $u$'s first interaction with a concept. An edge is
accepted when support exceeds 100 learners, $\mathrm{prec} \geq 0.75$ and
$\mathrm{lift} > 1.1$. Requiring lift is essential: precedence alone admits pairs
separated only by the calendar. The resulting graph is acyclic without intervention.

\subsection{Knowledge-Graph-Enhanced Knowledge Tracing}

Concept representations come from a three-layer graph convolution over the symmetrically
normalised adjacency of the prerequisite graph,
%
\begin{equation}
\hat{A} = D^{-1/2}(A + I)D^{-1/2},
\label{eq:adjacency}
\end{equation}
%
where $A \in \mathbb{R}^{C \times C}$ is the weighted adjacency over $C = 237$ concepts.
With $H^{(0)} = E$ a learned embedding table,
%
\begin{equation}
H^{(l+1)} = \mathrm{Drop}\!\left(\mathrm{ReLU}\!\left(\mathrm{LN}\!\left(\hat{A} H^{(l)} W^{(l)}\right)\right)\right), \quad l = 0,1,2,
\label{eq:gcn}
\end{equation}
%
giving $Z = H^{(3)} \in \mathbb{R}^{C \times 64}$. $Z$ is recomputed at every forward
pass, so gradients propagate into the graph structure.

For interaction $t$ with concept $c_t$, response $r_t$ and kind $k_t$,
%
\begin{equation}
x_t = W_c Z[c_t] + W_r e(r_t) + W_k e(k_t) + W_f f_t + W_s s,
\label{eq:embedding}
\end{equation}
%
where $s$ denotes static learner features and $f_t$ carries click volume, elapsed time
since the previous interaction, elapsed time since the same concept was last seen, and
position through the presentation.

Forgetting is modelled as an additive attention bias decaying with elapsed time,
%
\begin{equation}
b^{(h)}_{ij} = -\,\mathrm{softplus}(\gamma_h)\,\log(1 + \Delta_{ij}),
\label{eq:decay}
\end{equation}
%
with $\Delta_{ij}$ the days between interactions $i$ and $j$ and $\gamma_h$ learned per
head, so that one head may attend to recent activity while another retains longer
memory. Attention is
%
\begin{equation}
A^{(h)} = \mathrm{softmax}\!\left(\frac{Q^{(h)}{K^{(h)}}^{\top}}{\sqrt{d_k}} + b^{(h)} + M\right),
\label{eq:attention}
\end{equation}
%
where $M$ combines the causal and padding masks. Setting $\gamma_h = 0$ recovers ordinary
attention and constitutes the temporal ablation of Table~\ref{tab:ablation}.

Writing $h_t$ for the encoder state after interaction $t$,
%
\begin{equation}
\hat{y}_t = \sigma\!\left(\mathrm{MLP}\!\left([\,h_{t-1}\,;\,Z[c_t]\,]\right)\right),
\label{eq:prediction}
\end{equation}
%
conditioning on the identity of $c_t$ but never its response. Training minimises binary
cross-entropy over assessment positions $\mathcal{S}$ alone,
%
\begin{equation}
\mathcal{L} = -\frac{1}{|\mathcal{S}|}\sum_{t \in \mathcal{S}} \left[\, y_t \log \hat{y}_t + (1 - y_t)\log(1 - \hat{y}_t)\,\right].
\label{eq:loss}
\end{equation}
%
Clickstream interactions update the hidden state without contributing a loss term. This
dual-signal construction is what makes a sequence model of length 200 trainable on a
dataset in which learners average 8.7 assessments each.

Because predicted mastery is displayed to learners, probabilities are calibrated by
temperature scaling \cite{guo2017calibration},
%
\begin{equation}
T^\star = \operatorname*{arg\,min}_{T} \mathrm{BCE}\!\left(\sigma(z/T), y\right), \qquad \hat{y} = \sigma(z/T^\star),
\label{eq:temperature}
\end{equation}
%
fitted on one half of each fold's held-out predictions and evaluated on the other. The
mastery vector consumed by the planner evaluates the prediction head against every
concept,
%
\begin{equation}
m_t[c] = \sigma\!\left(\mathrm{MLP}\!\left([\,h_t\,;\,Z[c]\,]\right)\right), \quad c = 1,\dots,C.
\label{eq:mastery}
\end{equation}

\subsection{Prerequisite-Constrained Path Optimisation}

Recommendation is formalised as a Markov decision process. The state concatenates the
mastery vector, a coverage mask of concepts already recommended, the static learner
features and the remaining budget,
%
\begin{equation}
s_t = [\, m_t \,\|\, o_t \,\|\, s \,\|\, b_t \,] \in \mathbb{R}^{2C + 57}.
\label{eq:state}
\end{equation}
%
An action $a \in \{1, \dots, C\}$ is admissible only if
%
\begin{equation}
m_t[a] < \tau_{\mathrm{m}} \;\wedge\; \forall p \in \mathrm{pre}(a): m_t[p] \geq \tau_{\mathrm{r}} \;\wedge\; a \in \mathcal{E},
\label{eq:mask}
\end{equation}
%
with $\mathcal{E}$ the concepts of the learner's enrolled modules. Encoding the
pedagogical constraint in the admissible set, rather than expecting a learned policy to
discover it, makes an unprepared recommendation impossible by construction; this is
verified empirically in Table~\ref{tab:pathquality}. The reward
%
\begin{equation}
r_t = \alpha \sum_c \max(0, \Delta m[c]) + \beta \mathbf{1}[\mathrm{prereq}] - \delta \mathbf{1}[\mathrm{redundant}] - \eta
\label{eq:reward}
\end{equation}
%
uses $\alpha = 1.0$, $\beta = 0.2$, $\delta = 0.5$, $\eta = 0.05$.

Predicted mastery is bimodal in this setting: because non-submission is treated as
failure (Section~\ref{sec:implementation}), engaged learners score near $0.98$ on nearly
every concept and disengaged learners near $0.00$. Fixed thresholds therefore fail to
discriminate, and $\tau_{\mathrm{m}}, \tau_{\mathrm{r}}$ are set to the $70$th and $30$th
percentiles of each learner's own mastery distribution.

The policy evaluated here is greedy with respect to one-step predicted improvement,
%
\begin{equation}
\mathrm{score}(a) = \sum_{c \in \mathcal{E}} \left[\, \mathrm{logit}\, m^{(a)}_{t+1}[c] - \mathrm{logit}\, m^{(\varnothing)}_{t+1}[c] \,\right].
\label{eq:score}
\end{equation}
%
Two choices in \eqref{eq:score} are load-bearing. The baseline $m^{(\varnothing)}_{t+1}$
is an equal elapsed period \emph{without} study rather than the present state: studying
advances the clock, and the decay of \eqref{eq:decay} then attenuates every earlier
interaction, so that measuring against the present moment makes every available action
appear harmful. Scoring in log-odds rather than probability is necessary because engaged
learners saturate near $0.98$, where probability differences round to zero and the
ranking degenerates. A learned policy over \eqref{eq:state}--\eqref{eq:reward} is left to
future work; Section~\ref{sec:results} reports how much headroom remains for one.

\subsection{Explanation by Priority Backtracking}

For a recommended concept $c^\star$, each prerequisite ancestor $a$ is scored by its
shortfall, discounted by graph distance,
%
\begin{equation}
\pi(a) = \max(0,\ \tau_{\mathrm{r}} - m[a]) \cdot \rho^{\,d(a,c^\star) - 1},
\label{eq:priority}
\end{equation}
%
with $\rho = 0.6$ and traversal depth at most four, so that an unmet prerequisite
immediately preceding the recommendation outranks an equal shortfall further back.
Ancestors are reported in descending $\pi$. Counterfactual statements are computed by
raising a gap concept's mastery to $\tau_{\mathrm{m}}$ and re-evaluating
\eqref{eq:score}, so the claimed effect of closing a gap is measured rather than
asserted. Algorithm~\ref{alg:recommend} states the full procedure.

\begin{algorithm}[t]
\caption{Explained recommendation}
\label{alg:recommend}
\begin{algorithmic}[1]
\Require history $h$, graph $G$, parameters $\theta$, budget $k$
\Ensure $k$ recommendations, each with an explanation
\State $Z \gets \textsc{Gcn}(\hat{A})$ \Comment{Eq.~\eqref{eq:gcn}}
\State $m \gets \textsc{Mastery}(h, Z, \theta)$ \Comment{Eq.~\eqref{eq:mastery}, calibrated}
\State $\tau \gets \textsc{Percentiles}(m|_{\mathcal{E}})$
\State $L \gets \{a : \textsc{Admissible}(a, m, \tau, G)\}$ \Comment{Eq.~\eqref{eq:mask}}
\If{$L = \varnothing$} \State $L \gets \textsc{Relax}(m, \tau)$ \EndIf
\State $m^{(\varnothing)} \gets \textsc{Mastery}(\textsc{Idle}(h), Z, \theta)$
\ForAll{$a \in L$} \Comment{single batched pass}
  \State $m^{(a)} \gets \textsc{Mastery}(h \oplus (a, \mathrm{correct}), Z, \theta)$
  \State $\mathrm{score}(a) \gets \sum_{c \in \mathcal{E}}\!\left[\mathrm{logit}\,m^{(a)}[c] - \mathrm{logit}\,m^{(\varnothing)}[c]\right]$
\EndFor
\State $A \gets \textsc{TopK}(L, \mathrm{score}, k)$
\ForAll{$a \in A$}
  \State $\mathcal{G} \gets \textsc{Backtrack}(G, m, a, \tau)$ \Comment{Eq.~\eqref{eq:priority}}
  \State $\mathcal{C} \gets \textsc{Counterfactual}(a, \mathcal{G}, \theta)$
  \State $\textsc{Render}(a, \mathcal{G}, \mathcal{C})$
\EndFor
\State \Return $A$ with explanations
\end{algorithmic}
\end{algorithm}

% =====================================================================
\section{Implementation}
\label{sec:implementation}

\subsection{Environment}

The framework is implemented in Python 3.11 with PyTorch. Data preparation is performed
in SQL against an embedded DuckDB database, so that every transformation is auditable
without reading application code. Training used a single NVIDIA T4; inference and all
data processing run on CPU. Table~\ref{tab:efficiency} reports measured cost.

\subsection{Dataset and Preprocessing}

Evaluation uses the Open University Learning Analytics Dataset \cite{kuzilek2017oulad},
comprising 32{,}593 enrolments, 10{,}655{,}280 clickstream events and 173{,}912
assessment records across seven modules. Categorical attributes are one-hot encoded and
missing assessment scores are median-imputed. OULAD timestamps are day-granular with no
intra-day resolution; a session is therefore taken to be one learner-day.

Two properties of the data required specific treatment.

\textbf{Label construction.} OULAD records only \emph{submitted} assessments. Binarising
at the pass mark of 40 yields a 95.6\% positive rate, with the fifth percentile of
submitted scores equal to the pass mark itself; the failures are absent from the
records rather than from reality. Of 323{,}925 expected submissions, 150{,}013 never
occurred. Treating non-submission as failure, restricted to assessments falling while a
learner remained registered, yields a $67.9/32.1$ balance and changes the prediction
target to successful completion.

\textbf{Concept supervision is sparse and uneven.} Assessments exist for only 89 of the
237 concepts. Mastery estimates for the remaining 148 are reached through prerequisite
propagation, which is the condition Section~\ref{sec:results} examines directly.

Sequences are split $80/10/10$ and evaluated by five-fold cross-validation stratified on
final outcome and \emph{grouped by learner}, so that no learner contributes to both
training and test partitions. Table~\ref{tab:dataset} summarises the resulting corpus.

\subsection{Model Configuration}

The knowledge tracing encoder uses two transformer layers, hidden dimension 128, four
attention heads and a maximum sequence length of 200. The graph encoder uses three
convolution layers with layer normalisation and ReLU, producing 64-dimensional concept
embeddings. Dropout is $0.2$. Optimisation uses Adam at learning rate $10^{-3}$ with
cosine annealing, weight decay $10^{-4}$, batch size 64 and at most 12 epochs with early
stopping at patience 3. The model has 466{,}569 parameters. Random seeds are fixed and
recorded with every result.

\subsection{Evaluation Protocol}

Knowledge tracing is assessed by AUC-ROC, RMSE and expected calibration error over
assessment positions on held-out learners. Statistical comparisons use paired $t$-tests
over the five folds with Holm--Bonferroni correction across the comparison family, and
Cohen's $d$ is reported alongside, since with large numbers of paired decisions
statistical significance is attainable at negligible effect size.

Path quality is assessed against observed behaviour rather than simulation: each
held-out learner's history is truncated at several points, each planner is asked for a
recommendation, and the ground truth is the set of concepts that learner actually
engaged with over the following interactions. Only learners who passed or achieved
distinction contribute ground truth.

% =====================================================================
\section{Results and Discussion}
\label{sec:results}

\subsection{Predictive Accuracy}

Table~\ref{tab:accuracy} reports predictive performance. The proposed framework attains
AUC-ROC $0.9538 \pm 0.0019$, exceeding DKT by $0.0103$ ($p = 0.0009$) and the
graph-based baseline by $0.0036$ ($p = 0.022$) after correction. A majority-class
baseline returns exactly $0.5000$, confirming that the reported discrimination is not an
artefact of the $67.9/32.1$ class distribution. SAKT --- attention without graph
conditioning or temporal decay --- reaches $0.9233$, indicating that the improvement does
not follow simply from the use of attention.

Calibration is reported in Table~\ref{tab:calibration}. Because predicted mastery is
shown to learners inside explanations, a poorly calibrated probability is an incorrect
statement rather than a cosmetic defect; expected calibration error of $0.0054$ after
temperature scaling supports the values the interface displays.

\subsection{Why Conventional Ablation Is Insufficient}

Table~\ref{tab:ablation} reports the standard component ablation. Removing
knowledge-graph propagation changes accuracy by $-0.0001$ ($p = 0.77$) and removing
temporal attention by $-0.0006$ ($p = 0.45$). Read at face value, neither component
contributes.

That reading is an artefact of the measurement. Predictive accuracy is computed only at
assessment positions, and assessments exist for 89 of the 237 concepts, each carrying on
the order of $2{,}700$ labelled examples. A free embedding table learns such concepts
adequately without structural information; the graph is redundant there. Its function is
to support the remaining 148 concepts, which never enter the accuracy computation. The
conventional ablation is structurally incapable of observing the component it removes.

\subsection{Cold-Concept Evaluation}

To measure the contribution directly, all assessments were withheld for 27 of the 89
assessed concepts and their observed responses removed from the model input, so that
they behave as genuinely unassessed concepts; evaluation was then restricted to those
concepts alone. Table~\ref{tab:coldconcept} reports the outcome: $0.9004 \pm 0.0059$
with graph propagation against $0.8679 \pm 0.0042$ without, a difference of $+0.0325$
($p = 0.0003$, $d = 5.06$), with the graph-conditioned model superior on every fold.

Knowledge-graph propagation therefore contributes nothing where direct supervision is
plentiful and substantially where it is absent. This has a practical consequence for the
recommender, which ranks across all 237 concepts, most of which carry no assessment; and
a methodological consequence for the literature, in which component ablations are
routinely reported on metrics unable to observe the component under test.

Table~\ref{tab:coldstart} reports the related case of learners rather than concepts. With
no interaction history the model reaches $0.7421$ using a learner-profile prior obtained
by clustering early behaviour, rising to $0.9449$ after four interactions.

\subsection{Learning Path Quality}

Table~\ref{tab:pathquality} reports 906 recommendation decisions from 400 held-out
successful learners. The proposed planner attains the highest NDCG@5 at $0.1532$, ahead
of weakest-first ($0.1344$), random selection among admissible actions ($0.1300$),
popularity ($0.1069$) and syllabus order ($0.0612$). All differences are significant
after correction.

The effect sizes qualify that conclusion and we report them rather than the $p$-values
alone. Against syllabus order the effect is small ($d = 0.36$); against random selection
among admissible actions it is negligible ($d = 0.09$). The honest interpretation is that
the prerequisite mask performs most of the work and the learned ranking adds a small but
consistent improvement. This also bounds the headroom available to a learned policy: an
agent optimising over \eqref{eq:reward} must improve on a greedy baseline that itself
exceeds random selection by $d = 0.09$.

Prerequisite violations are zero across all 906 decisions and every planner. The
pedagogical guarantee of \eqref{eq:mask} is therefore verified rather than assumed.
Concept coverage of $0.726$ indicates that the planner does not collapse onto a small
set of popular concepts, the characteristic failure mode that popularity-based
recommendation exhibits here at $0.510$.

\subsection{Explanation Quality}

Table~\ref{tab:explanation} reports objective measures over 25 recommendation decisions,
each repeated under three perturbations. The sample is small and we report it as such;
the measures are nonetheless informative because each is computed against the model's own
counterfactual response rather than against an opinion. Fidelity of $0.5318$ indicates
that cited prerequisite deficits are positively but imperfectly predictive of measured
counterfactual improvement. Stability of $0.5310$ indicates that roughly half the cited
gaps persist under a $\pm 0.02$ perturbation of the mastery estimate.

Sufficiency of $0.0020$ is the most informative of the three and we report it as a
negative finding. It shows that the concept named in an explanation accounted for only
$0.2\%$ of the predicted benefit, the remainder arising from transfer to related concepts
the explanation did not mention. The explanations were accurate but substantially
incomplete. The renderer was subsequently amended to name the principal beneficiaries.
No rating scale would have surfaced this: participants cannot assess a quantity the
explanation omits.

% =====================================================================
\section{Conclusion and Future Work}
\label{sec:conclusion}

\subsection{Conclusion}

This paper presented an explainable learning-path recommendation framework combining
knowledge-graph-conditioned knowledge tracing, prerequisite-constrained planning and
priority-backtracking explanation, evaluated on OULAD. The knowledge tracing model
attains AUC-ROC $0.9538 \pm 0.0019$ with expected calibration error $0.0054$,
significantly exceeding DKT, SAKT and a graph-based baseline.

The principal finding concerns evaluation. A conventional ablation indicates that
knowledge-graph propagation contributes nothing, yet under a protocol that evaluates
concepts lacking direct supervision the same component improves accuracy by $+0.0325$
($d = 5.06$). The discrepancy is not a property of this model but of the metric: aggregate
accuracy measured where supervision is abundant cannot observe a structural prior whose
purpose is to compensate for its absence. We suggest that component ablations in this
literature be accompanied by evidence that the reporting metric is capable of detecting
the component under test.

\subsection{Limitations}

The prediction target is successful completion rather than knowledge. Because OULAD
records only submitted assessments, treating non-submission as failure is necessary for
a usable class balance but means a substantial share of the signal is engagement rather
than understanding. This is appropriate for a system intended to prompt intervention,
but it is a weaker claim than measuring mastery.

The concept graph is derived, not given. OULAD provides no concept annotations,
prerequisite relations or material titles, so concepts are defined structurally and
labelled by activity composition rather than subject matter. Placement for 82\% of
materials relies on observed access; the substitute is validated at $\rho = 0.975$ on the
subset carrying both signals, but deriving concept order from behaviour and then mining
prerequisites from behaviour is partially self-referential.

Supervision covers 89 of 237 concepts, so mastery for the remainder is not directly
calibrated. Predicted mastery is bimodal, requiring learner-relative thresholds. The
proposed model additionally conditions on learner features that the DKT and SAKT
baselines do not, as those follow their original formulations; part of the reported
improvement may therefore reflect feature access rather than architecture.

All results derive from a single institution's distance-learning data across seven
modules, and generalisation to campus-based teaching is untested.

\subsection{Ethical Considerations}

All data used in this study is drawn from OULAD, which is publicly released under
CC-BY 4.0 and contains no personally identifying information; learners are represented by
anonymised identifiers. No data was collected from human participants for the experiments
reported here. The system produces advisory recommendations and makes no determinative
decision concerning any learner.

The model conditions on attributes including gender, disability status, index of multiple
deprivation and age band. A subgroup analysis of predictive accuracy and recommendation
quality across these attributes has not been conducted and is identified below as
required future work. Until it is, the system should not be deployed in any setting where
its recommendations influence assessment or progression decisions.

\subsection{Future Work}

Four directions follow directly from the limitations above.

\textbf{Fairness analysis.} Subgroup performance across gender, disability, deprivation
band and age band must be established before deployment. This is the most immediate
requirement.

\textbf{A learned policy.} Section~\ref{sec:results} quantifies the available headroom:
the greedy planner exceeds random selection among admissible actions by $d = 0.09$. A
policy learned over \eqref{eq:state}--\eqref{eq:reward}, evaluated both in simulation and
by off-policy estimation on logged trajectories, would establish whether that margin can
be widened.

\textbf{User evaluation.} The objective measures reported here establish whether
explanations are faithful to the model; they do not establish whether learners find them
useful. A controlled study with ethical approval, comparing conditions with and without
explanation, would address the complementary question.

\textbf{Expert validation of mined edges.} The prerequisite graph is inferred from
behaviour. Rating a sample of mined edges by subject experts would provide an independent
estimate of its precision.

% =====================================================================
\section*{Reproducibility}

OULAD is available at \texttt{research.stem.open.ac.uk/ouanalyse/dataset} under CC-BY 4.0.
All source code, SQL transformations, configurations and generated results are available
at \texttt{github.com/abhay27singh/explainable-\allowbreak learning-path-recommender}.
Every table in this paper is regenerated from raw data by a single command with fixed
random seeds.

% =====================================================================
%                              TABLES
% =====================================================================

\begin{table*}[t]
\centering
\caption{Comparative predictive performance. Mean $\pm$ standard deviation over five-fold
cross-validation with folds grouped by learner. $p$-values are paired $t$-tests against
the proposed model, Holm--Bonferroni corrected; $d$ is Cohen's $d$ for paired samples.}
\label{tab:accuracy}
\begin{tabular}{lccccc}
\toprule
Model & AUC-ROC & RMSE & ECE & $p$ & $d$ \\
\midrule
Majority (no-skill)         & $0.5000 \pm 0.0000$ & $0.4669 \pm 0.0004$ & $0.0037$ & $<0.001$ & --- \\
SAKT \cite{pandey2019sakt}  & $0.9233 \pm 0.0018$ & $0.2992 \pm 0.0017$ & $0.0081$ & $<0.001$ & $26.16$ \\
DKT \cite{piech2015dkt}     & $0.9435 \pm 0.0017$ & $0.2751 \pm 0.0022$ & $0.0102$ & $0.0009$ & $5.67$ \\
GNN-based                   & $0.9502 \pm 0.0014$ & $0.2665 \pm 0.0020$ & $0.0064$ & $0.0220$ & $2.25$ \\
\textbf{Proposed}           & $\mathbf{0.9538 \pm 0.0019}$ & $\mathbf{0.2621 \pm 0.0026}$ & $\mathbf{0.0054}$ & --- & --- \\
\midrule
Improvement vs.\ GNN-based  & $+0.0036$ & $-0.0044$ & & & \\
Improvement vs.\ DKT        & $+0.0103$ & $-0.0130$ & & & \\
\bottomrule
\end{tabular}
\end{table*}

\begin{table}[t]
\centering
\caption{Component ablation under the conventional protocol. Neither component produces
a detectable change. Section~\ref{sec:results}-B explains why this measurement cannot
observe the knowledge graph.}
\label{tab:ablation}
\begin{tabular}{lccc}
\toprule
Variant & AUC-ROC & RMSE & $p$ \\
\midrule
Proposed (full)              & $0.9538 \pm 0.0019$ & $0.2621 \pm 0.0026$ & --- \\
\quad w/o knowledge graph    & $0.9539 \pm 0.0021$ & $0.2618 \pm 0.0025$ & $0.77$ \\
\quad w/o temporal attention & $0.9532 \pm 0.0014$ & $0.2624 \pm 0.0024$ & $0.45$ \\
\bottomrule
\end{tabular}
\end{table}

\begin{table}[t]
\centering
\caption{Cold-concept ablation. All assessments are withheld for 27 of the 89 assessed
concepts and their responses removed from the model input; evaluation is restricted to
those concepts. The graph-conditioned model is superior on every fold.}
\label{tab:coldconcept}
\begin{tabular}{lcc}
\toprule
Variant & AUC-ROC & RMSE \\
\midrule
With knowledge graph & $\mathbf{0.9004 \pm 0.0059}$ & $\mathbf{0.3359 \pm 0.0059}$ \\
Without              & $0.8679 \pm 0.0042$ & $0.3611 \pm 0.0113$ \\
\midrule
Difference & \multicolumn{2}{c}{$+0.0325$, \ $p = 0.0003$, \ $d = 5.06$} \\
\bottomrule
\end{tabular}
\end{table}

\begin{table*}[t]
\centering
\caption{Learning path quality over 906 recommendation decisions from 400 held-out
learners who passed or achieved distinction. Ground truth is the set of concepts each
learner actually engaged with next. All planners operate under the same admissible-action
constraint of \eqref{eq:mask}.}
\label{tab:pathquality}
\begin{tabular}{lccccc}
\toprule
Planner & NDCG@5 & Hit@5 & Prec@5 & Prereq.\ violations & Coverage \\
\midrule
Syllabus order      & $0.0612 \pm 0.154$ & $0.1832$ & $0.0453$ & $0.0000$ & $0.452$ \\
Popularity          & $0.1069 \pm 0.199$ & $0.2936$ & $0.0808$ & $0.0000$ & $0.510$ \\
Random (admissible) & $0.1300 \pm 0.218$ & $0.3455$ & $0.0863$ & $0.0000$ & $0.722$ \\
Weakest-first       & $0.1344 \pm 0.207$ & $0.3642$ & $0.0956$ & $0.0000$ & $0.515$ \\
\textbf{Proposed}   & $\mathbf{0.1532 \pm 0.238}$ & $\mathbf{0.3687}$ & $\mathbf{0.1051}$ & $\mathbf{0.0000}$ & $\mathbf{0.726}$ \\
\bottomrule
\end{tabular}

\vspace{3pt}
{\footnotesize
\begin{tabular}{lccc}
\multicolumn{4}{c}{Paired comparison of NDCG@5 against the proposed planner (Holm-corrected)}\\
\toprule
Comparison & Difference & $p$ & Cohen's $d$ \\
\midrule
vs.\ syllabus order  & $+0.0921$ & $<0.001$ & $0.36$ (small) \\
vs.\ popularity      & $+0.0464$ & $<0.001$ & $0.17$ (negligible) \\
vs.\ random          & $+0.0232$ & $0.0125$ & $0.09$ (negligible) \\
vs.\ weakest-first   & $+0.0189$ & $0.0167$ & $0.08$ (negligible) \\
\bottomrule
\end{tabular}}
\end{table*}

\begin{table}[t]
\centering
\caption{Calibration of predicted mastery. Temperature is fitted on one half of each
fold's held-out predictions and evaluated on the other.}
\label{tab:calibration}
\begin{tabular}{lc}
\toprule
& Expected calibration error \\
\midrule
Before temperature scaling & $0.0084$ \\
After temperature scaling  & $\mathbf{0.0054}$ \\
\bottomrule
\end{tabular}
\end{table}

\begin{table}[t]
\centering
\caption{Predictive accuracy as a function of available interaction history.}
\label{tab:coldstart}
\begin{tabular}{lc}
\toprule
Prior interactions & AUC-ROC \\
\midrule
$0$        & $0.7421 \pm 0.0353$ \\
$1$--$4$   & $0.9449 \pm 0.0086$ \\
$5$--$9$   & $0.9559 \pm 0.0062$ \\
$10$--$19$ & $0.9539 \pm 0.0056$ \\
$\geq 20$  & $0.9514 \pm 0.0023$ \\
\bottomrule
\end{tabular}
\end{table}

\begin{table}[t]
\centering
\caption{Objective explanation quality over 25 recommendation decisions, each with
three perturbation replicates. These measures record whether an explanation is faithful
to the model, rather than whether a reader found it agreeable.}
\label{tab:explanation}
\begin{tabular}{lcp{3.1cm}}
\toprule
Measure & Value & Definition \\
\midrule
Fidelity    & $0.5318$ & Rank correlation between cited deficits and measured counterfactual improvement \\
Sufficiency & $0.0020$ & Share of predicted gain attributable to the concept named \\
Stability   & $0.5310$ & Overlap of cited gaps under $\pm 0.02$ perturbation \\
\bottomrule
\end{tabular}
\end{table}

\begin{table}[t]
\centering
\caption{Dataset and derived concept graph.}
\label{tab:dataset}
\begin{tabular}{lr}
\toprule
Quantity & Value \\
\midrule
Distinct learners               & $28{,}785$ \\
Enrolments                      & $32{,}593$ \\
Clickstream events              & $10{,}655{,}280$ \\
Assessment records              & $173{,}912$ \\
Learning materials              & $6{,}364$ \\
\midrule
Concepts derived                & $237$ \\
\quad prerequisite edges        & $724$ \\
\quad corequisite edges         & $302$ \\
\quad concepts with assessments & $89$ of $237$ \\
\midrule
Interaction sequences           & $27{,}904$ \\
Supervised tokens               & $242{,}670$ \\
Label balance                   & $67.9\% / 32.1\%$ \\
Placement validation            & $\rho = 0.975$ \\
\bottomrule
\end{tabular}
\end{table}

\begin{table}[t]
\centering
\caption{Computational cost, measured on the hardware used.}
\label{tab:efficiency}
\begin{tabular}{lr}
\toprule
Quantity & Value \\
\midrule
Model parameters                  & $466{,}569$ \\
Recommendation latency (CPU)      & ${\approx}160$\,ms \\
Cached mastery lookup             & $7$\,ms \\
Training per configuration (T4)   & ${\approx}25$\,min \\
Data preparation from raw files   & ${\approx}2.5$\,min \\
\bottomrule
\end{tabular}
\end{table}

% =====================================================================
\begin{thebibliography}{00}

\bibitem{li2026survey}
A.~Li, Y.~Li, and X.~Gao, ``Personalized learning path recommendation based on knowledge
graphs: A survey,'' \emph{Electronics}, vol.~15, no.~1, Art. no. 238, 2026,
doi: 10.3390/electronics15010238.

\bibitem{aher2013combination}
S.~B. Aher and L.~M.~R.~J. Lobo, ``Combination of machine learning algorithms for
recommendation of courses in e-learning system based on historical data,''
\emph{Knowledge-Based Systems}, vol.~51, pp.~1--14, 2013,
doi: 10.1016/j.knosys.2013.04.015.

\bibitem{kuzilek2017oulad}
J.~Kuzilek, M.~Hlosta, and Z.~Zdrahal, ``Open University Learning Analytics dataset,''
\emph{Scientific Data}, vol.~4, Art. no. 170171, 2017, doi: 10.1038/sdata.2017.171.

\bibitem{corbett1994knowledge}
A.~T. Corbett and J.~R. Anderson, ``Knowledge tracing: Modeling the acquisition of
procedural knowledge,'' \emph{User Modeling and User-Adapted Interaction}, vol.~4,
no.~4, pp.~253--278, 1994.

\bibitem{piech2015dkt}
C.~Piech, J.~Bassen, J.~Huang, S.~Ganguli, M.~Sahami, L.~Guibas, and J.~Sohl-Dickstein,
``Deep knowledge tracing,'' in \emph{Advances in Neural Information Processing Systems},
vol.~28, 2015, pp.~505--513.

\bibitem{zhang2017dkvmn}
J.~Zhang, X.~Shi, I.~King, and D.-Y. Yeung, ``Dynamic key-value memory networks for
knowledge tracing,'' in \emph{Proc. 26th Int. Conf. World Wide Web (WWW)}, 2017,
pp.~765--774, doi: 10.1145/3038912.3052580.

\bibitem{ghosh2020akt}
A.~Ghosh, N.~T. Heffernan, and A.~S. Lan, ``Context-aware attentive knowledge tracing,''
in \emph{Proc. 26th ACM SIGKDD Int. Conf. Knowledge Discovery and Data Mining (KDD)},
2020, pp.~2330--2339, doi: 10.1145/3394486.3403282.

\bibitem{pandey2019sakt}
S.~Pandey and G.~Karypis, ``A self-attentive model for knowledge tracing,'' in
\emph{Proc. 12th Int. Conf. Educational Data Mining (EDM)}, 2019, pp.~384--389.

\bibitem{vaswani2017attention}
A.~Vaswani, N.~Shazeer, N.~Parmar, J.~Uszkoreit, L.~Jones, A.~N. Gomez, L.~Kaiser, and
I.~Polosukhin, ``Attention is all you need,'' in \emph{Advances in Neural Information
Processing Systems}, vol.~30, 2017, pp.~5998--6008.

\bibitem{guo2017calibration}
C.~Guo, G.~Pleiss, Y.~Sun, and K.~Q. Weinberger, ``On calibration of modern neural
networks,'' in \emph{Proc. 34th Int. Conf. Machine Learning (ICML)}, 2017,
pp.~1321--1330.

\bibitem{holm1979}
S.~Holm, ``A simple sequentially rejective multiple test procedure,''
\emph{Scandinavian Journal of Statistics}, vol.~6, no.~2, pp.~65--70, 1979.

\bibitem{ma2024survey}
B.~Ma, T.~Yang, and B.~Ren, ``A survey on explainable course recommendation systems,''
in \emph{Distributed, Ambient and Pervasive Interactions}, N.~A. Streitz and S.~Konomi,
Eds. Cham: Springer, 2024, pp.~273--287, doi: 10.1007/978-3-031-60012-8\_17.

\end{thebibliography}

\end{document}
